\section{Nested Benchmarks}\label{sec:benchmarks}

The contribution depends on an interaction, so the appropriate comparison is not whether each component is individually familiar. It is whether the mechanism in Theorem~\ref{thm:main} remains available after either strategic margin is removed. The benchmarks below are nested interventions on the full game, not additional global-equilibrium theorems.

In the \emph{fixed-primary-location} benchmark, the primary-location profile is assigned exogenously before the policy comparison and does not respond to any policy deviation; the backup game, disaster-state probabilities, and product-market continuation are otherwise solved exactly as in the baseline model. In the \emph{fixed-private-redundancy} benchmark, readiness levels are treated as exogenous constants during the policy comparison and do not respond to protection; the location stage, primary-risk mapping, and product-market continuation remain active. Thus each benchmark removes exactly one endogenous response while retaining the other baseline primitives. Table~\ref{tab:benchmarks} summarizes the logic.

\begin{table}[t]
\centering
\caption{Nested benchmark logic}
\label{tab:benchmarks}
\small
\begin{tabular}{p{.20\textwidth}p{.25\textwidth}p{.46\textwidth}}
\toprule
Model & Endogenous margins & Policy consequence \\
\midrule
Fixed primary locations & Private redundancy only & Public/private substitution remains, but protection cannot attract primary production. \\
Fixed private redundancy & Primary location only & Standard local public-input competition remains, but public policy cannot reorganize private resilience. \\
Full model & Primary location and private redundancy & Plant attraction operates after redundancy crowd-out changes the social value of protection; $a^{LG}>0=a^{SP}$ can arise. \\
\bottomrule
\end{tabular}
\end{table}

\subsection{Fixed primary locations}

Fix any primary-location profile and suppress the location stage. A unilateral protection increase then cannot attract a primary plant: the expected local plant count is constant and the derivative of the $2bp$ term in \eqref{eq:government-payoff} with respect to local protection is zero. The backup game is still re-solved after every policy change, so public protection can continue to crowd out private readiness. This benchmark therefore isolates public/private resilience substitution without the plant-attraction channel.

\subsection{Fixed private redundancy}

Fix the firms' readiness levels exogenously and hold them unchanged under every policy deviation. Primary locations are still chosen through the Bayesian location game, and protection still changes primary disruption risk and location attractiveness. The model therefore retains competition for mobile production through a local public input, but public policy can no longer reorganize firms' contingency capacity. This benchmark isolates the plant-attraction channel without endogenous resilience adjustment.

\subsection{Full-game interaction}

The full model combines the two margins but does more than place them side by side. Public protection changes the private resilience architecture, which changes the national value of additional protection. Local governments then act on a different objective because they value the induced relocation of primary production. The canonical calculations verify the direction of each benchmark mechanism, while Theorem~\ref{thm:main} and Appendix~\ref{app:verification} certify the global policy-ranking reversal only for the full game. The benchmark comparison is therefore a mechanism-identification exercise, not a claim that each restricted game separately satisfies the full theorem's global policy ranking.
